\section{Recovery of the physical low estimate}
\label{gex:low-recovery}

The additive polynomial \(A_{m,\sigma,\nu}(Y')\) is the one in the
exact separation \cite[(6.2)]{OAI26}.  Its fixed ray coefficient is
\(\mathbf1_{s\in\sigma}\chi_s(b_*)/(\tau\xi(s))\), extended by zero
off the original primary support before evaluating the inverse.
For \(Q,Y'\ge1\), in fixed polynomial ranges, and
\begin{equation}
 P_a=Y'^2/Q\ge1,
\label{gex:additive-gram-parameter}
\end{equation}
\cite[Proposition~15.2]{OAI26} gives
\begin{equation}
 \sum_{q_m\ll Q}|A_{m,\sigma,\nu}(Y')|^2
 \ll_{\varepsilon,\mathrm{fixed}}
 (1+|\nu|)^{J_A}\frac Q{Y'}
 \left(1+P_a^{1/6}+\frac{P_a^2}{Y'}\right)Z^\varepsilon.
\label{gex:additive-gram-input}
\end{equation}
The order \(J_A\) is independent of the moving scales and height.
The same input covers fixed finite combinations of these ray coefficients.

At \eqref{gex:rescaled-low-scales},
\begin{equation}
 P_a=q_{b_*}^{-1}Z^b,\qquad b=1/8.
\label{gex:gram-rescaling-invariant}
\end{equation}
For \(\varpi=1/6000\), every \(0\le d\le\ell\) satisfies
\begin{equation}
\begin{aligned}
 M'&\ge\tfrac12-3\varpi>0,\\
 l_x-d&\ge\tfrac3{16}-\tfrac32\varpi>0,\\
 l_y-d-11b/6&\ge\tfrac1{12}-\tfrac32\varpi>0.
\end{aligned}
\label{gex:uniform-low-scale-margins}
\end{equation}
The bounded tuple ratios therefore give \(Q,Y',P_a\ge1\) eventually,
uniformly in \(J\), and \(Y'\gg P_a^{11/6}\).  Both the constant
term and \(P_a^2/Y'\) in \eqref{gex:additive-gram-input} are then
\(O(P_a^{1/6})\).  Hence
\begin{equation}
 \sum_{q_m\ll Q}|A_{m,\sigma,\nu}(Y')|^2
 \ll (1+|\nu|)^{J_A}(Q/Y')P_a^{1/6}Z^\varepsilon.
\label{gex:additive-gram-at-new-scales}
\end{equation}

The following proposition is \cite[Proposition~15.3]{OAI26} at the
perturbed geometry; its proof repeats that proof, with
Lemma~\ref{gex:completed-low} in place of \cite[Lemma~15.1]{OAI26}.
At \(\varpi=0\) the bound is \cite[(15.6)]{OAI26}.

\begin{proposition}[cf.\ {\cite[Proposition~15.3]{OAI26}}]
\label{gex:physical-low}
For the geometry \eqref{gex:geometry} and every \(\varepsilon>0\),
the complete compensated probe satisfies
\begin{equation}
 |\mathcal I_{\eta,\mathrm{mod}}(Z)|
 \ll_{\varepsilon,\mathrm{fixed}}
 Z^{l_x/2+b/12+\varepsilon}
 =Z^{3/16-1/24000+\varepsilon}.
\label{gex:physical-low-bound}
\end{equation}
The datum, windows and slot count are fixed before \(Z\);
all subsets and tuples in \eqref{gex:compensated-source} are included.
\end{proposition}

\begin{proof}
For a fixed rescaled tuple, apply Cauchy in the shared physical rows
of \cite[(6.2)]{OAI26}, whose normalization is \(Q^{-1/2}\).
Lemma~\ref{gex:completed-low} and
\eqref{gex:additive-gram-at-new-scales} give
\begin{equation}
 Q^{-1/2}
 \left(\frac Q{Y'}P_a^{1/6}\right)^{1/2}
 Z^{M'/2+(5\ell-1+d)_+/8+\varepsilon}
 \ll (X')^{1/2}P_a^{1/12}
       Z^{(5\ell-1+d)_+/8+\varepsilon}.
\label{gex:fixed-tuple-low}
\end{equation}
We used \(Z^{M'}/(X'Y')=r_J^2\), a bounded factor.
The polynomial height costs from the two inputs are integrable
against the common, rapidly decreasing Mellin density.  Choosing
its finite derivative order after the tail orders absorbs those
costs uniformly in the tuple.

Ideal counting gives \(O(Z^{d+\varepsilon})\) tuples on \(J\).
Their coefficient is \(q_{p_J}^{-3/2}\asymp Z^{-3d/2}\), and
\((X')^{1/2}\asymp X^{1/2}Z^{-d/2}\).  Their total additional
exponent in \eqref{gex:fixed-tuple-low} is
\begin{equation}
 d-3d/2-d/2+(5\ell-1+d)_+/8
 \le-7d/8\le0,
\label{gex:full-tuple-payment}
\end{equation}
because \(\ell<1/5\).  The surviving \(J^c\) weights, including
\(\overline{\eta(p_{J^c})}\), were already summed inside
\(B_{m,\sigma}^J\).

Sum the finitely many subsets and fixed ray sectors, allocating
their small-power losses inside \(\varepsilon\).
The Gaussian and Mellin integrations have the finite weighted
norms used above.  Finally,
\(P_a^{1/12}\asymp Z^{b/12}\) and
\(l_x/2+b/12=3/16-\varpi/4\), proving the proposition.
\end{proof}

Since \(C(s)=s-11/16\) in \eqref{gex:principal-exponent},
\(C(20999/24000)=3/16-1/24000\).
Thus the full low estimate matches the boundary of the principal signal.
